\documentclass[10pt, letterpaper]{article}
\usepackage[utf8]{inputenc}
\usepackage[T1]{fontenc}
\usepackage{geometry}
\usepackage{enumitem}
\usepackage{hyperref}
\usepackage{titlesec}
\usepackage{color}
\usepackage{fancyhdr}

% Margins
\geometry{top=0.75in, bottom=0.75in, left=1in, right=1in}

% Colors
\definecolor{columbiablue_official}{RGB}{0, 33, 71}

% Header
\pagestyle{fancy}
\fancyhf{}
\renewcommand{\headrulewidth}{0.5pt}
\renewcommand{\headrule}{\hbox to\headwidth{\color{columbiablue_official}\leaders\hrule height \headrulewidth\hfill}}
\lhead{\small \textbf{\color{columbiablue_official} COLUMBIA UNIVERSITY} \\ \href{https://www.arch.columbia.edu/}{Graduate School of Architecture, Planning and Preservation}}
\rhead{\small Protocol Summary \\ \href{https://github.com/dhardestylewis/thesis/tree/main/Submitted/IRB_Submitted/Protocol_and_Submission}{ACYY0820}}

% Section Styling
\titleformat{\section}{\large\bfseries\color{columbiablue_official}}{}{0em}{}[\color{columbiablue_official}\titlerule]
\titlespacing*{\section}{0pt}{10pt}{4pt}

\begin{document}

\begin{center}
    {\Large \textbf{\color{columbiablue_official} Study Protocol Summary}} \\[0.3em]
    {\large Predicting NIMBYism in Austin’s Housing Reform Era}
\end{center}

\section*{Key Personnel}
\noindent \textbf{Principal Investigator (PI):} Hiba Bou Akar, Ph.D. \\
Associate Professor, Urban Planning \\
Columbia University GSAPP \\
Email: \href{mailto:hb2541@columbia.edu}{hb2541@columbia.edu}

\vspace{0.5em}
\noindent \textbf{Student Researcher:} Daniel Lewis \\
M.S. Urban Planning Student, Columbia University GSAPP \\
Email: \href{mailto:dl3645@columbia.edu}{dl3645@columbia.edu}

\section*{IRB Contact}
If you have questions about your rights or welfare as a research participant, you may contact: \\
\textbf{Columbia University Morningside Institutional Review Board (IRB)} \\
Phone: (212) 851-7040 \\
Email: \href{mailto:AskIRB@columbia.edu}{AskIRB@columbia.edu}

\section*{Study Purpose}
Austin, Texas faces a severe housing affordability crisis driven by rapid population growth, tech-sector expansion, and historically restrictive land-use regulations. Recent housing reforms have coincided with organized neighborhood opposition to development using tools such as protest petitions, public testimony, and legal challenges. This study aims to:
\begin{enumerate}[leftmargin=*, noitemsep]
    \item Empirically characterize patterns of neighborhood opposition to housing development using administrative records (2018–2025).
    \item Examine the democratic implications of predictive tools that could forecast opposition, through interviews with key stakeholders.
\end{enumerate}

\section*{Research Questions}
\textbf{Primary:} How can cities identify patterns of neighborhood opposition, and what are the democratic implications of using predictive analytics to anticipate it? \\
\textbf{Sub-questions:}
\begin{itemize}[leftmargin=*, noitemsep]
    \item What demographic and geographic factors predict opposition to specific projects?
    \item How do stakeholders perceive the legitimacy of algorithmic tools in housing policy?
    \item What alternative technological approaches could improve processes without surveillance concerns?
\end{itemize}

\section*{Study Design}
This mixed-methods, minimal-risk study combines:
\begin{itemize}[leftmargin=*, noitemsep]
    \item \textbf{Record Review:} Analysis of public records (zoning cases, protest petitions, campaign contributions, census data). De-identified datasets will use coded IDs.
    \item \textbf{Interviews:} Semi-structured interviews (45–60 mins) with ~10–15 stakeholders (city officials, neighborhood leaders, advocates, developers).
    \item \textbf{Observation:} Review of public meetings (Council, commissions).
\end{itemize}

\section*{Participation Details}
\textbf{Interviews:} Conducted in-person or via secure videoconference. With permission, they are audio-recorded for transcription. Participation is voluntary; you may decline to answer or stop at any time.

\noindent \textbf{Risks:} Minimal risk. Potential breach of confidentiality is mitigated by secure storage, encryption, and reporting only aggregate/de-identified results.

\noindent \textbf{Benefits:} No direct benefit. Indirect benefits include informing more equitable planning practices.

\end{document}
